%========================= INTRODUÇÃO =========================
\section{Introdução}

% RASCUNHO. A estrutura e o fio do argumento estão prontos; falta
% inserir as citações, que só podem entrar depois da revisão.
% Onde está [CITAR], entra a referência do artigo correspondente.

Fobias específicas estão entre as condições de maior prevalência
relacionadas à ansiedade, e as associadas a animais figuram entre as mais
comuns delas [CITAR]. Cobras, ratos, baratas, aranhas e sapos aparecem de
forma recorrente nessa categoria --- animais que, no contexto urbano,
raramente oferecem risco objetivo, mas cuja simples possibilidade de
encontro basta para alterar a rotina de quem os teme.

A limitação que daí decorre é menos visível que a reação em si. Ela se
manifesta em evitação: deixar de entrar num cômodo, de visitar um lugar,
de exercer uma atividade. [CITAR]

\subsection*{O impasse}

O procedimento reconhecido para a superação desse tipo de medo envolve o
contato progressivo e controlado com o estímulo temido [CITAR]. E é
justamente aí que se instala um impasse: \emph{o contato é precisamente
aquilo que a pessoa evita}. O procedimento depende daquilo que a
condição impede.

Ambientes virtuais oferecem uma via para contornar esse impasse. Neles, o
estímulo pode ser graduado com precisão --- aproximado, multiplicado,
posto em movimento --- enquanto o participante mantém a certeza de que
nada ali é real e de que pode interromper a qualquer instante. O controle
permanece com quem é exposto, o que é condição do procedimento e não um
detalhe de conforto.

\subsection*{A lacuna}

% COMPLETAR com o que a revisão efetivamente mostrar. A redação abaixo
% é a hipótese de partida e precisa ser confirmada ou corrigida.
A literatura que investiga exposição a estímulos fóbicos em Realidade
Virtual concentra-se, em larga medida, em estudos conduzidos em
laboratório, com equipamento dedicado e acompanhamento presencial
[CITAR]. São condições que asseguram controle experimental, mas que
restringem o alcance: dependem de um local específico, de um aparelho
específico e de uma pessoa presente para operá-lo.

A consolidação do WebXR --- especificação que permite executar conteúdo
tridimensional imersivo diretamente no navegador --- altera esse quadro,
ao viabilizar que a mesma aplicação funcione em computador, celular e
óculos sem instalação. Discute-se menos, contudo, o que essa mudança
acarreta: como as pessoas percebem e interagem com estímulos fóbicos
quando o ambiente virtual é alcançado por um endereço, e não por um
laboratório.

\subsection*{Questão de pesquisa e objetivo}

Deste trabalho, portanto, a questão:

\begin{quote}
Como as pessoas percebem e interagem com estímulos associados a fobias de
animais em ambientes virtuais organizados em níveis progressivos de
aproximação?
\end{quote}

Para investigá-la, foi desenvolvido um protótipo que simula cinco animais
associados a fobias, cada um no ambiente em que seria encontrado, com a
exposição organizada em cinco níveis nos quais variam a distância, a
quantidade de exemplares, o movimento e a possibilidade de interação. A
aplicação registra, durante o uso, medidas de comportamento que são
confrontadas com o grau de medo que o participante declara previamente.

\subsection*{Delimitação de escopo}

Cabe esclarecer, desde a introdução, que o sistema desenvolvido é um
\textbf{protótipo acadêmico}. Ele não constitui tratamento psicológico ou
médico, não realiza diagnóstico e não se propõe substituto de
acompanhamento profissional. A avaliação conduzida restringe-se a
usabilidade, imersão, interação e percepção do estímulo.

Essa delimitação não é ressalva protocolar: ela determina o método. É por
ela que não se aplica instrumento clínico de rastreio, que não há
critério de inclusão por diagnóstico e que o grau de medo é tratado como
autorrelato de percepção, e não como medida clínica.

\subsection*{Organização do texto}

A Seção~2 apresenta os objetivos. A Seção~3 reúne o referencial teórico
construído a partir da revisão da literatura. A Seção~4 descreve a
metodologia, incluindo o protocolo de busca, o perfil dos participantes e
as variáveis coletadas. A Seção~5 descreve o artefato desenvolvido e as
decisões de projeto que o determinaram. A Seção~6 apresenta e discute os
resultados. A Seção~7 conclui, separando o que os dados sustentam do que
permanece como conjectura.
